\section{递归定义与结构归纳}
\label{sec:01-Mathematical-Language/04-Induction-and-Recursion/03}

你写了一个函数，把一份 JSON 文档里所有数值字段加起来。对象里可以套数组，数组里可以再套对象，嵌套深度没有上限。几十份样本都跑通了，但``任意深''这三个字，样本一份也没触及。

想把``对所有合法文档都正确''说清楚，先得说清楚什么叫合法文档。而合法文档也有无穷多份。上一节的强归纳依赖``所有更小的''，可两份 JSON 谁比谁小，并不是天生就有答案。

\subsection{用有限条规则钉死一个无限集合}

先试着定义``算术表达式''：变量和常数是表达式；两个表达式用 \(+\) 或 \(\times\) 连起来还是表达式。听着像是把话说完了，但定义里用上了正在被定义的词。这种自指不会引出矛盾，前提是规则写全。

\begin{definition}
一个递归定义由三部分组成：若干\textbf{基础对象}，不借助别的对象就被判定为合法；若干\textbf{构造规则}，从已经合法的对象造出新的合法对象；以及一条\textbf{封闭条款}，声明除前两条能造出来的以外，再没有别的成员。
\end{definition}

封闭条款常被省略不写，可它才是定义的边界。有了它，``某个古怪的东西算不算表达式''就有了裁决依据：不算，因为造不出来。更要紧的是，它保证每个成员都有一条长度有限的生成历史——这正是下面那套证明方法的立足点。

熟悉的数列定义是同一模板的最简形态。阶乘写成 \(0!=1\) 与 \((n+1)!=(n+1)\,n!\)，第一条是基础对象，第二条把 \((n+1)!\) 的计算责任转交给 \(n!\)。只写第二条会怎样？计算一路向更小的下标退行，永远落不了地。Fibonacci 需要两个基础值，因为它的构造规则回望两步。

数据类型也照搬这个模板。一棵二叉树要么是单个叶子，要么由两棵已有的二叉树挂在一个新根下面得到，此外别无二叉树。JSON 文档的定义同理：数字、字符串、布尔值和 null 是基础对象；由已有文档组成的数组、以及字段值为已有文档的对象，是构造规则。

顺带澄清一件容易混的事。递归定义说的是对象\emph{是什么}，不是在推荐实现方式。Fibonacci 的定义照抄成代码会得到一个正确但极慢的程序，因为它重复计算了海量子问题。数学定义、正确算法、高效算法是三个可以分开讨论的层面。

\subsection{证明跟着构造语法走}

对递归生成的对象，证明的形状应当和定义的形状一致。结构归纳只有两步：先验证所有基础对象满足性质；再对每一条构造规则，假设它的各个输入对象都满足性质，证明它的输出也满足。

\begin{center}
\begin{tikzpicture}[x=1cm,y=1cm,>=stealth,font=\small]
  \node[draw,circle,inner sep=2.5pt] (r)  at (0,2.4) {\(\times\)};
  \node[draw,circle,inner sep=2.5pt] (p)  at (-1.5,1.2) {\(+\)};
  \node[draw,circle,inner sep=2.5pt] (x1) at (1.5,1.2) {\(x\)};
  \node[draw,circle,inner sep=2.5pt] (x2) at (-2.4,0) {\(x\)};
  \node[draw,circle,inner sep=2.5pt] (c3) at (-0.6,0) {\(3\)};
  \draw (r) -- (p); \draw (r) -- (x1); \draw (p) -- (x2); \draw (p) -- (c3);
  \draw[->,thick,black!55] (2.9,0) -- (2.9,2.4);
  \node[right,align=left,black!70] at (3.05,1.2) {结构归纳\\沿这个方向推进};
  \node[below,black!60] at (-1.5,-0.45) {基础对象：变量与常数};
\end{tikzpicture}

\emph{表达式 \((x+3)\times x\) 的构造历史就是这棵树；性质从叶子出发，每经过一条构造规则被验证一次，抵达根节点时已覆盖整个表达式。}
\end{center}

为什么这样就够了？因为封闭条款保证每个合法对象都有有限的生成历史。任取一个对象，沿它的生成历史回溯，基础步管住了起点，归纳步管住了每一次构造，有限步之后必然合拢。上一节要求的良基顺序在这里也有了着落：``更小''就是``在生成历史中更早''，生成步数是一个取值于自然数、严格下降的度量。

于是不必先把树、字符串、表达式编码成自然数再套普通归纳。沿构造语法一层层走上去就行，而且每条规则只需处理一次。

\subsection{一个短示范：满二叉树的叶子与内部节点}

命题：任意有限满二叉树（每个内部节点恰好两个子节点）中，叶子数 \(L\) 与内部节点数 \(I\) 满足 \(L=I+1\)。

基础对象是单个叶子：\(L=1\)，\(I=0\)，等式成立。

构造步：取满足等式的两棵树 \(T_1,T_2\)，挂到新根下得到 \(T\)。新根贡献一个内部节点，叶子数是两棵子树之和，所以 \(L=L_1+L_2\)、\(I=I_1+I_2+1\)，代入得

\[
L=(I_1+1)+(I_2+1)=(I_1+I_2+1)+1=I+1 .
\]

两步都通过，而每棵有限满二叉树恰由这两条规则生成，结论对全体成立。

整个过程不必关心树长什么形状，不必按层数节点，也不必额外对高度做一次归纳。开头那个 JSON 求和函数的正确性论证是同一副骨架：基础对象上函数返回该数值本身或零，数组与对象两条构造规则上，函数返回各子文档结果之和——逐条对上，就覆盖了任意深度。

\subsection{同一个对象，两条构造路径}

递归定义要能支撑证明，还有一个前提容易被跳过：由构造路径定义出来的东西，不能依赖于选了哪条路径。

用分数 \(p/q\) 表示有理数时，若规定某个量为 \(p+q\)，麻烦立刻出现：\(1/2\) 与 \(2/4\) 是同一个数，却给出 3 和 6 两个答案。这时需要单独证明良定义性，即结果与代表元的选择无关。

表达式那边的对策是让构造路径唯一。裸写 \(a+b\times c\) 有两种读法，保留括号或直接以解析树为对象，就使每个表达式恰好对应一棵树，歧义消失。写下递归公式只是第一步，函数是否良好定义是另一道独立的检查。

结构本身也可能悄悄超出定义。一张计算图里，某个中间结果若被下游多处引用，它就不再是树而是有向无环图；递归遍历仍然可行，但必须加备忘，否则同一个节点会被重复计入。倘若连``沿依赖走不会回到起点''都不成立——出现了环——递归定义就失去了基础情形，整套方法要重来。对象的生成结构变了，跟着它走的证明也得跟着变。

树造完就摆在那里，静止不动。下一节换成动态的东西：循环的状态一轮一轮生成，生成规则是``执行一次循环体''，而要证的性质得在整条执行轨迹上一直成立。
